\section{The Coming}

\subsection{Go Out to Meet Him}\label{sec:meet}

On the First Sunday of Advent the priest opens the Mass of the new year with a
prayer for going out. He asks almighty God that his faithful may go out with
just works to meet the coming Christ, be placed at his right hand, and possess
the kingdom of heaven. The Church's year begins on the road.

The prayer is new in its place and old in its matter. The reform of the Missal
recast it from a prayer of the old Roman book called the Gelasian Sacramentary,
where it was said after Communion in an Advent Mass. That older prayer asked
that God's whole household be given this will:
\latin{Domino nostro venienti in operibus iustis aptos occurrere}, ``to run to
meet our Lord as he comes, made ready in just works.'' And it went on to the
judgment: \latin{et eius dexterae sociati regnum mereantur possidere caeleste},
``and that, joined to his right hand, they may merit to possess the heavenly
kingdom.'' Other early books used the same prayer for daily Masses and as a
prayer over the people, and one Gallican book set it as the blessing of the
people at the first Mass of Advent. The prayer that now opens the year was once
said over a congregation about to leave the church.

To go out to meet him is the cry of the Lord's own parable of the ten virgins:

\begin{quote}
And at midnight there was a cry made: Behold the bridegroom cometh. Go ye forth
to meet him.\par
{\smallskip\noindent\small ---Matthew 25:6\par}
\end{quote}

\noindent The Latin of the cry is \latin{exite obviam ei}. The Office of the
same Sunday hears it again. In the Office of Readings, after St Cyril of
Jerusalem has preached the two comings of Christ, the Church sings a
responsory, a chant after the reading in which a refrain returns between the
verses. It is built on the one the Roman Breviary sang at Matins, the night
Office whose place the Office of Readings now takes. The Breviary's form ran:

\begin{quote}
I look from afar, and, behold, I see the Power of God coming, and a cloud
covering all the land. Go ye out to meet Him, and say: Tell us if Thou art He,
That shalt reign over God's people Israël.\par
{\smallskip\noindent\small ---Responsory \latin{Aspiciens a longe}, the Roman Breviary, in Bute's English\par}
\end{quote}

\noindent Its refrain in Latin is \latin{Ite obviam ei, et dicite}, ``Go out to
meet him, and say,'' and the present Office keeps its call to go out to meet
him.

What does it mean to go out? St Augustine put the question to his people and
gave them a plain reply:

\begin{quote}
They wished to go to meet the Bridegroom. What is the meaning of ``to go and
meet the Bridegroom''? To go with the heart, to be waiting for his coming.\par
{\smallskip\noindent\small ---St Augustine, Sermon 93, on the Ten Virgins\par}
\end{quote}

\noindent But the prayer of the First Sunday speaks
also of just works, and the parable of the virgins turns on oil. St John
Chrysostom, preaching on the same parable, said what the oil was and
where it is bought:

\begin{quote}
But by lamps here, He meaneth the gift itself of virginity, the purity of
holiness; and by oil, humanity, almsgiving, succor to them that are in need.
\ldots{} ``But go to them that sell, and buy.'' And who are they that sell? The
poor. And where are these? Here, and then should they have sought them, not at
that time. \ldots{} Seest thou what great profit arises to us from the poor?
shouldest thou take them away, thou wouldest take away the great hope of our
salvation. \ldots{} Wherefore here must we get together the oil, that it may be
useful to us there, when the time calls us. For that is not the time of
collecting it, but this.\par
{\smallskip\noindent\small ---St John Chrysostom, Homily 78 on St Matthew\par}
\end{quote}

\noindent St Thomas Aquinas linked the Fathers' comments on the Gospels verse
by verse into a single chain, the \emph{Catena Aurea}, or Golden Chain. St
Hilary, as the Catena Aurea gives him, says that the poor, in return for the
alms of the faithful, sell them the consciousness of good works, and he calls
this ``the abundant fuel of an undying light which may be bought and stored up
for the fruits of mercy.''

The right hand of the old prayer is the right hand of the Last Judgment, where
the King sorts the nations:

\begin{quote}
Then shall the king say to them that shall be on his right hand: Come, ye
blessed of my Father, possess you the kingdom prepared for you from the
foundation of the world.\par
{\smallskip\noindent\small ---Matthew 25:34\par}
\end{quote}

\noindent The sentence goes on to the hungry fed and the naked clothed, and the
King counts it all as done to himself. Chrysostom, preaching on that page,
weighed every reason the King gives, and among them this: ``The dignity of the
one receiving, for it was God, who was receiving by the poor.''

A week later, on the Second Sunday, the Collect takes up the same road. It
speaks of those who hurry out to meet the Son, and it too reworks a prayer of
the Gelasian book, which began with the same haste.

And the season ends where the first Sunday points. When Christmas came, St
Gregory Nazianzen began his sermon on the Lord's birth with
a summons the Church has been rehearsing since the end of November:

\begin{quote}
Christ is born, glorify ye Him. Christ from heaven, go ye out to meet Him.
Christ on earth; be ye exalted. Sing unto the Lord all the whole earth; and
that I may join both in one word, Let the heavens rejoice, and let the earth be
glad, for Him Who is of heaven and then of earth.\par
{\smallskip\noindent\small ---St Gregory Nazianzen, Oration 38, On the Birthday of Christ\par}
\end{quote}
